\begin{lemma}
\label{lemma-ample-after-field-extension}
Let $k$ be a field. Let $X$ be a scheme over $k$.
If there exists an ample invertible sheaf on $X_K$ for some
field extension $K/k$, then $X$ has an ample invertible
sheaf.
\end{lemma}

\begin{proof}
Let $K/k$ be a field extension such that $X_K$ has
an ample invertible sheaf $\mathcal{L}$.
The morphism $X_K \to X$ is surjective. Hence $X$ is quasi-compact
as the image of a quasi-compact scheme (Properties, Definition
\ref{properties-definition-ample}). Since $X_K$ is quasi-separated
(by Properties, Lemma
\ref{properties-lemma-affine-s-opens-cover-quasi-separated})
we see that $X$ is quasi-separated: If $U, V \subset X$ are
affine open, then $(U \cap V)_K = U_K \cap V_K$ is quasi-compact
and $(U \cap V)_K \to U \cap V$ is surjective. Thus
Schemes, Lemma \ref{schemes-lemma-characterize-quasi-separated} applies.

\medskip\noindent
Write $K = \colim A_i$ as the colimit of the subalgebras of $K$
which are of finite type over $k$. Denote
$X_i = X \times_{\Spec(k)} \Spec(A_i)$.
Since $X_K = \lim X_i$ we find an $i$ and an invertible sheaf
$\mathcal{L}_i$ on $X_i$ whose pullback to $X_K$ is $\mathcal{L}$
(Limits, Lemma \ref{limits-lemma-descend-invertible-modules};
here and below we use that $X$ is quasi-compact and quasi-separated as
just shown). By Limits, Lemma \ref{limits-lemma-limit-ample}
we may assume $\mathcal{L}_i$ is ample after possibly increasing $i$.
Fix such an $i$ and let $\mathfrak m \subset A_i$ be a maximal
ideal. By the Hilbert Nullstellensatz
(Algebra, Theorem \ref{algebra-theorem-nullstellensatz})
the residue field $k' = A_i/\mathfrak m$ is a finite
extension of $k$. Hence $X_{k'} \subset X_i$ is a closed subscheme
hence has an ample invertible sheaf
(Properties, Lemma \ref{properties-lemma-ample-on-closed}).
Since $X_{k'} \to X$ is finite locally free we conclude
that $X$ has an ample invertible sheaf by
Divisors, Proposition \ref{divisors-proposition-push-down-ample}.
\end{proof}